% 修订正文：保留六个承担建模任务的公式；数值已展开；与独立修订稿采用同一证据。
\section{问题四的模型与求解}
\label{sec:q4r}

\subsection{定向接收与观测约束}

第三问以测角反馈缩小源的位置区域，并协调搜索、定位与清除。加入定向源后，几何上有利的测点可能位于辐射背面，无法获得方位信息；无信号也同时具有超距和背向两种解释。若仍沿用全向源的排除规则，就可能删去真实位置。因此，本问构造适用于任意朝向的搜索布局，并结合有信号与无信号的历史观测，联合约束位置、半径和朝向，形成\textbf{联合域规划策略}。

Tokekar 与 Isler\cite{Tokekar2013}用有界测角误差构造几何可行域，为前文的扇形求交提供了依据。Vander Hook 等\cite{VanderHook2014}同时考虑移动与测量时间，沿主要不确定方向的垂线谨慎接近目标，可用于建立测点选择基线。其高斯估计与方位歧义设定不同于本题，应用时需改写状态更新方式，并另行处理定向源的接收条件。

针对方向性，赵倩倩等\cite{Zhao2025}在无人机部署中考虑目标发射天线的方向，提示应将源位置与朝向共同作为未知量；与其连续增益模型不同，本文采用题设的半圆可见性条件。Song 等\cite{Song2012}依据多个未知无线电源的空间概率分布引导机器人搜索，提供了利用更新后的位置分布规划运动的思路。其 RSS 似然依赖信号强度，本题仅反馈方位和接收状态，故下文重新构造相容条件。

沿用第三问的区域 $\Omega$、位置域 $K$、最小包围圆中心 $c(K)$ 和半径 $\rho(K)$。以 $x$ 表示未知源位置 $g$ 的候选值，作用半径 $R\in[1000,1500]$ m，类型 $z=0,1$ 分别表示全向与定向，单位向量 $n$ 表示发射朝向。机器人在 $q$ 处的可接收状态为
\begin{equation}
 h(q;x,R,z,n)=\mathbf1\{\|q-x\|\le R\}
 \mathbf1\{z=0\ \text{或}\ n^{\mathsf T}(q-x)\ge0\}.
 \label{eq:q4r-observation}
\end{equation}
有效接收后还要满足测角误差或 5 m 近场条件。这里的发射朝向不同于测得的源方位，图\ref{fig:q4r-direction}说明了这种差异。

\begin{figure}[htbp]\centering
\includegraphics[width=\linewidth]{fig01-direction-coverage.pdf}
\caption{无信号歧义与方向鲁棒搜索原理。左图中背向和超距均产生无信号；右图中源位于站点凸包内，任意朝向下至少一站位于前向。此图为机制示意。}
\label{fig:q4r-direction}\end{figure}

\Needspace{9\baselineskip}
\subsection{从谨慎测向到联合域规划}

沿用第三问的文献适配思路，构造\textbf{谨慎测向策略}。以位置域凸包的均匀面积均值$\mu$和协方差近似原文估计，令$\sigma_{\max}$为最大特征值的平方根，$v$为相应特征向量的垂直方向。按文献\cite{VanderHook2014}的 Algorithm 1，在$\mu\pm av$中选择离当前位置较近者，其中
\begin{equation}
 a=\frac{\sigma_{\max}}{\sqrt{\sigma_{\beta_{\mathrm c}}^2-\sigma_s^2}},\qquad
 \sigma_{\beta_{\mathrm c}}=\frac{\pi}{2\Phi^{-1}(1-\beta_{\mathrm c}/2)}.
 \label{eq:q4r-cautious}
\end{equation}
本问在测试前固定$\beta_{\mathrm c}=0.05$，$\sigma_s^2$取$\pm1^\circ$均匀误差的二阶矩。位置域仍由有界几何约束更新，清除仍按包围圆判定。若无信号使区域不变，且两个候选点都已测量，则以$30^\circ$步长旋转同半径候选；无可用测点时转入有限光学覆盖。该策略保留原文的谨慎选点思想，以适应本题的固定测角误差和定向接收条件。

谨慎测向只比较交会几何和移动距离，未评价候选点的可接收性。为此，先构造\textbf{可见性评分策略}：以包围圆中心代表目标，用有限的半径、朝向假设估计接收可能性，将失联后的恢复时间纳入测点评分，并共同安排目标与搜索站点的访问次序。这样可避开部分低效测点，但无信号仅影响行动评分，没有直接缩小连续位置域；包围圆中心也未区分候选位置与观测历史的相容程度。

\textbf{联合域规划策略}进一步将位置、半径、类型与朝向共同代入历史观测，删除已不可能成立的位置块，再按相容程度计算用于路线规划的代表位置。表\ref{tab:q4r-models}概括了三种策略的差别。三者均使用22站搜索布局、路线组合、顺路测量、数量上界停止规则和有限清除回退。各策略代码与参数在相应演练前固定，比较重点是接收历史如何参与选点和位置更新。

\begin{table}[htbp]\centering\small
\caption{对照策略与逐步解决的问题}\label{tab:q4r-models}
\begin{tabularx}{\linewidth}{L{2.3cm} Y Y}\toprule
策略 & 决策依据 & 仍存在的问题或新增作用\\\midrule
谨慎测向 & 沿主要不确定方向的垂线测量，选较近侧 & 有利交会角未必可接收，可能反复失联后进入光学覆盖\\
简洁可见性评分 & 有限假设估计可见性与行动费用，包围圆中心用于规划 & 无信号不直接裁剪连续位置域，目标代表位置较粗\\
联合域规划 & 保留 可见性评分 框架，加入联合域证书和后验代表位置 & 将接收历史同时用于几何排除与运动规划；保证条件与启发式排序分开\\\bottomrule
\end{tabularx}\end{table}

\subsection{任意朝向下的搜索覆盖}

仅以接收圆盘覆盖目标区域不足以保证定向源被发现。对任意候选源位置 $x$，若存在站点子集 $S_x$ 满足
\begin{equation}
 x\in\operatorname{conv}(S_x),\qquad
 \max_{s\in S_x}\|s-x\|\le1000\ \mathrm m,
 \label{eq:q4r-cover}
\end{equation}
则任意朝向下均能在至少一站收到信号。证明只需将 $x$ 写成站点的凸组合：各站相对 $x$ 的有向投影的加权和为零，因此不可能所有正权重站点均严格位于背向；至少一站前向且不超出最小作用半径。

采用中心、半径 970 m 的七点等角内环及半径 1850 m 的十四点等角外环，共 22 个站点，两环均从正横轴起始。外环略超出源分布圆域，为边界源提供外侧观测位置。对覆盖半径 1800 m 圆域的 4,988 个矩形单元，逐一检查四角位于见证站点凸包内，且到各见证站点距离均小于 1000 m；最大见证距离为 999.975 m。矩形的凸性与距离函数的凸性将四角检查推广到整个单元，因而这是一份连续区域证书，而非离散源位置的抽样检验。该布局是可行配置，不主张站点数最少。

每站扫描尚未发现的频道，并保存包括无信号在内的全部反馈。访问完全部站点后，覆盖证书保证搜索完整；若提前发现 16 个不同源，则由题目数量上界提前结束搜索。随后继续处理已发现未清除源，直至全部完成。搜索结束由覆盖完成或达到数量上界判定。

\subsection{无信号约束下的联合区域裁剪}

对每个频道，将与全部观测历史相容的参数组成联合域$\Theta_t$，并以$K_t$保守外包其位置投影。有效方位和近场反馈仍按第三问分别与扇形、5 m圆盘求交；无信号反馈则须结合此前接收位置，判断超距与背向能否由同一组参数解释。

设有效接收位置集合为 $\mathcal S^+$、无信号位置集合为 $\mathcal S^-$。考虑一个由圆盘 $B(c,r)$ 外包的位置块。有效接收给出整块共享的半径下界，并确定无法由超距解释的无信号点：
\begin{equation}
 L_B=\max\{1000,\max_{p\in\mathcal S^+}(\|p-c\|-r)\},\qquad
 \mathcal S^-_B=\{q\in\mathcal S^-:\|q-c\|+r<L_B\}.
 \label{eq:q4r-forced}
\end{equation}
若 $L_B>1500$，该块不可能包含源。若 $\mathcal S^-_B$ 非空，该块内任何相容源都必须为定向型，且对 $\mathcal S^-_B$ 中各点背向，同时对 $\mathcal S^+$ 中各点前向。这些限制共享同一个朝向，不能为每次观测分别任意选择解释。

对有效接收点取 $d=p-c$，对被迫背向点取 $d=c-q$。当 $\|d\|>r$ 时，块内位置变化引起的向量方向偏移至多为 $\arcsin(r/\|d\|)$；因此相容朝向必在以 $\arg d$ 为中心、半宽为 $\pi/2+\arcsin(r/\|d\|)$ 的闭圆周区间 $A(d,r)$ 内。当 $\|d\|\le r$ 时不加入该项限制。若 $\mathcal S^-_B\ne\varnothing$ 且
\begin{equation}
 \bigcap_{p\in\mathcal S^+}A(p-c,r)\ \cap\!
 \bigcap_{q\in \mathcal S^-_B}A(c-q,r)=\varnothing,
 \label{eq:q4r-reject}
\end{equation}
则整个位置块均可排除。由于各角区间由必要条件放宽得到，只有交集为空才能删除；非空或退化为单点的角区间均保留。这样便可利用无信号反馈，检验全部观测是否与同一接收半径和朝向相容。

实现沿最小旋转外接矩形的长边二分，块尺度降至 10 m 或处理数达到 4,096 时停止；未证伪区域全部保留。矩形和圆盘留微米量级外扩，距离判据留 $10^{-5}$ m 余量，角区间再放宽 $10^{-9}$ rad，避免边界误删。图\ref{fig:q4r-domain}给出实际日志中的裁剪。所示包围圆半径由 22.42 m 降至 17.40 m，低于 19.5 m 的无需补测清除阈值，因此可移至包围圆圆心完成清除。

\begin{figure}[htbp]\centering
\includegraphics[width=\linewidth]{fig02-joint-domain.pdf}
\caption{联合约束如何促成可靠清除。左图取全部 65 次半径严格下降事件按相对降幅排序的第 33 次，橙色斜线区为排除部分；右图展示 联合域规划 30 场模拟器演练测试的全部 113 次裁剪，65 次降低包围圆半径，半径比中位数为 0.9972。事件非独立场景，不将 113 当作性能样本量。}
\label{fig:q4r-domain}\end{figure}

\FloatBarrier
\subsection{路线、测点与清除位置的选择}

路线规划还需要一个能够代表目标位置的点。联合域规划在区域内取均匀位置、均匀半径与朝向，并对全向、定向类型赋予相同先验权重。对每个候选位置，将相容的半径和朝向积分后赋权，以加权平均位置生成路线节点和测点。位置积分用128个确定性 Sobol 面积样本近似；权重不足时退回几何代表点。上述先验服务于行动比较，不影响连续位置域的保守排除。

将剩余搜索站点和未清除目标的代表位置组成节点集，规划从当前位置出发、终点自由的开放路线。分别由最近邻、上一轮路线更新和最远种子插入构造初始解，经2-opt改进后取最短者。每完成一次站点扫描或一个目标的处理即重新规划。与第三问相同，本问也联合排序搜索和目标任务，但以代表位置近似目标处理位置，路线长度仅用于比较访问次序。

专用测点围绕代表位置及接近路径生成，包括位置域主轴的侧向点和中间点，剔除距已有测点不超过 2 m 的重复点。用至多 16 个代表位置、三个半径、24 个朝向及全向类型组成的历史相容假设，估计接收比例 $\widehat p(q)$ 和有信号后的平均包围圆半径。选择预计耗时最小的点：
\begin{equation}
 \widehat J(q)=t_{\mathrm{move}}(q)+t_{\mathrm{radio}}+
 t_{\mathrm{approach}}(q)+\widehat p(q)t_{\mathrm{loc}}(q)
 +(1-\widehat p(q))t_{\mathrm{recover}}(q).
 \label{eq:q4r-decision}
\end{equation}
各项依次为移动、一次测量、靠近目标、接收后继续定位及失联恢复的时间估计。移动速度和测量费用取题设值；继续定位时间由剩余半径及其相对清除尺度的对数估计，失联恢复时间由固定基础项和当前区域尺度估计，具体系数见配套代码。式\eqref{eq:q4r-decision}是候选点排序的启发式指标，其中有限假设等权；前述代表位置则由连续先验积分加权，两者承担不同计算任务。

到达搜索站点时，还对仍需定位、距历史测点足够远且视角变化或接近程度足够大的已知频道顺路测量。由此一次移动可以服务搜索与定位两类需求。若 $\rho(K_t)\le19.5$ m，则移至 $c(K_t)$ 清除，20 m 光学范围覆盖整个位置域；源的发射类型不影响这一判断。每次目标服务至多进行 10 次专用无线电测量，随后用边长 25 m 的相交方格中心实施有限光学覆盖，因为方格外接圆半径仅 17.68 m。失败后保守排除已检查圆盘，在未排除候选中继续选择。有限覆盖的完整性来自几何覆盖，顺序则按单位移动及失败费用的排除面积排序。

搜索布局与保守位置域分别提供发现和清除的几何依据，代表位置、路线和测点评分则用于降低耗时。有限光学覆盖可在测向进展缓慢时继续完成清除，但整场能否满足官方时间预算仍取决于实际行动费用。

\subsection{策略比较与结果分析}

\subsubsection{模拟器演练测试比较}

谨慎测向新增30次模拟器演练测试，可见性评分和联合域规划各采用既有30场完整记录，共90个不同案例。每场以总虚拟时间除以官方源数，再对各场结果取算术平均；均值区间采用29个自由度的 Student $t$分布，策略差异采用独立样本 Welch 区间。全部案例均纳入统计，结果见表\ref{tab:q4r-official}和图\ref{fig:q4r-official}。

\begin{table}[htbp]\centering\small
\caption{三种策略的模拟器演练测试比较}\label{tab:q4r-official}
\begin{tabularx}{\linewidth}{L{2.6cm} Z Z Z Z}\toprule
策略 & 全清场数 & 均值 & 均值 95\% 区间 & 第 90 百分位\\
 & & \multicolumn{3}{c}{时间单位：s/源}\\\midrule
谨慎测向 & 30/30 & 651.83 & [587.29, 716.36] & 825.11\\
可见性评分 & 30/30 & 496.70 & [456.42, 536.98] & 636.08\\
联合域规划 & 30/30 & \textbf{458.37} & [416.05, 500.69] & 588.52\\\bottomrule
\end{tabularx}\end{table}

\begin{figure}[htbp]\centering
\includegraphics[width=\linewidth]{fig03-official-comparison.pdf}
\caption{全部 90 场独立模拟器演练测试及分项计费。左图展示每场结果、场均值与 95\% 均值区间；右图先对每场分项时间除以该场源数再平均。光学费用包括失败尝试，图中仅统计虚拟时间。}
\label{fig:q4r-official}\end{figure}

联合域规划在30场中清除390个源，平均耗时458.37 s/源。相对谨慎测向，平均节省193.46 s/源（29.68\%），节省量95\%置信区间为$[117.67,269.25]$ s/源；相对可见性评分，平均节省38.33 s/源（7.72\%），区间为$[-18.85,95.51]$ s/源。前一差异较明确，后一差异仍可能由场景波动引起。

从已发现频道的后续观测看，谨慎测向出现617次无信号测量和928次失败光学尝试，可见性评分分别为393次、592次，联合域规划为194次、0次。谨慎测向的侧向点虽有利于交会，却可能落在辐射背面；可见性评分减少这类选择，联合域规划再利用历史反馈更新位置区域和规划代表点。

分项费用与上述变化一致：移动时间由谨慎测向的509.27降至联合域规划的333.31 s/源，占平均总时间差的约91\%；相对可见性评分，移动与无线电测量分别减少25.05和8.07 s/源。三组平均源数为12.77、12.87、13.00，定向源比例也不完全相同，因此这些费用和频次主要用于分析整体行为变化，各改进环节的独立贡献还需同场景消融实验。

\subsubsection{同场景路线比较}

在统一校准本地模拟器中，另固定30组场景及同点误差，对三种策略逐场比较。三者均完成全部清除，谨慎测向、可见性评分和联合域规划的均值分别为601.51、505.08、484.70 s/源。联合域规划较可见性评分在25场更快、5场更慢，平均节省20.38 s/源，配对差值95\%置信区间为$[11.38,29.37]$ s/源。本地场景和误差分布是对官方环境的近似，这组实验用于检查相同行动条件下的改进机制。

\begin{figure}[htbp]\centering
\includegraphics[width=\linewidth]{fig04-paired-paths.pdf}
\caption{三策略在同一场景的路径与清除进度。将全部 30 个本地场景按可见性评分与联合域规划的耗时差（前者减后者，单位为秒/源）排序，取第 15 场，含 10 个源，其中 3 个为定向源；圈为起点，叉为成功清除时的停留位置，不表示源的真实位置。坐标等比例且范围相同；进度曲线延续至搜索完成，末端方块表示任务结束，时间计入全部付费动作。}
\label{fig:q4r-paired}\end{figure}
\FloatBarrier

图\ref{fig:q4r-paired}所示场景中，可见性评分与联合域规划分别移动20145.21 m和19346.71 m，测量281次和273次，换频266次和268次。联合域规划少走798.49 m，节省159.70 s；少测8次，节省40 s；多换频2次，增加2 s，合计节省197.70 s。谨慎测向总耗时6112.14 s，联合域规划为5552.34 s。清除进度曲线表明，联合域规划并非始终领先，但较短的整体行程使其最终更早完成任务。

\Needspace{7\baselineskip}
联合域规划的官方日志中共排除1374个位置块，113次裁剪有65次降低包围圆半径，半径比中位数为0.9972。图\ref{fig:q4r-domain}说明裁剪在部分状态下可直接促成清除，但多数变化较小。另以仅采用相同后验代表位置、不进行联合裁剪的\textbf{后验规划策略}作历史对照，两者独立模拟器演练测试的均值差为12.53 s/源，95\%置信区间$[-41.84,66.90]$跨零。因此，当前实验尚不足以单独确定联合裁剪的时间收益。

\Needspace{18\baselineskip}
\subsection{正式测验结果}

联合域规划已完成三次官方正式测验，结果见表\ref{tab:q4r-formal}，官方历史页面均显示已上传。平均定位清除时间取累计虚拟时间与成功清除数之比；程序运行时间取官方接受退出和进入请求的时间戳之差，包含接口交互时间。

\begin{table}[htbp]\centering\small
\caption{联合域规划的三次官方正式测验结果}\label{tab:q4r-formal}
\begin{tabularx}{\linewidth}{L{5.6cm} Z Z Z}\toprule
案例编码 & 成功清除数 & 平均定位清除时间/s & 程序运行时间/s\\\midrule
9N77-TH84-ME6J-NRK4 & 10 & 557.293 & 4.325\\
QGEZ-KRWA-MNGB-6CMH & 11 & 525.395 & 4.426\\
KTZ8-EZQ5-937B-QNEV & 11 & 501.464 & 4.530\\\bottomrule
\end{tabularx}\end{table}

\FloatBarrier
